% =========================================================================
% 一、问题重述
% =========================================================================

\section{问题重述}

\subsection{问题背景}

微网是为解决光伏发电等分布式能源本地消纳而发展起来的小型发用电系统。出于经济性考虑，非孤岛式微网可在外部电网（外网）电价较低或分布式能源有剩余电量时为储能设备充电，在电价较高时释放电能。储能设备最大容量 12000 kWh，允许储电量范围 1200--10800 kWh，最大充放电功率 5000 kW，充放电效率均为 90\%；2025 年 1 月 1 日 0:00 储电量为 6000 kWh。微网供电不得低于小区负载，富余供能可弃置，不向外部电网售电。附件 1 给出代表日 10 分钟间隔的电价、小区负载与光伏发电预测功率；附件 2 给出 2025 年逐日实测负载与光伏功率。

\subsection{四问对比}

题面四问共享同一物理系统，差别只在决策时可用的信息与可调整的决策。表~\ref{tab:restate} 按``可用信息—可调决策—新增难点''三条线把四问压缩对照：可用信息决定决策时能看多远，可调决策刻画各问授予的自由度，新增难点则是相对前一问必须额外处理的部分。

\begin{table}[H]
  \centering
  \caption{四问的信息条件、决策自由度与核心新增难点}
  \label{tab:restate}
  \small
  \begin{tabularx}{\textwidth}{@{}l >{\raggedright\arraybackslash}X >{\raggedright\arraybackslash}X >{\raggedright\arraybackslash}X@{}}
    \toprule
    问题 & 决策时可用信息 & 可调整的决策 & 核心新增难点 \\
    \midrule
    Q1 & 给定代表日电价、负载及光伏输入 & 0:00 购电计划、确定性储能轨迹 & 损耗与峰谷价差 \\
    \addlinespace
    Q2 & 已发生历史与已知分时电价；当日未来实测不可用（本文信息假设） & 日前计划锁定，执行按当前净负载反馈 & 预测偏差与 5 倍紧急费 \\
    \addlinespace
    Q3 & 增加附件 3 多次发布的光伏预报 & 6/12/18 时可改未来合同 & 信息价值、调整成本、实际 SOC \\
    \addlinespace
    Q4 & 增加实时波动价格；未来价格不可见（本文解释） & 分别继承 Q2/Q3 权限 & 相对价格与跨日库存配置 \\
    \bottomrule
  \end{tabularx}
\end{table}

\subsection{必答交付物}

四问均要求给出可复算的具体结果。问题一至问题三按题目表 1 格式给出指定时间段的购电量、全天购电量与购电费，按表 2 格式给出指定时间段的充放电量与 0:00、24:00 储电量；问题二、问题三另按表 3 格式给出紧急购电结果。其中问题一针对代表日，\textbf{问题二、问题三还需分别给出 2025 年 3 月 20 日、6 月 21 日、9 月 23 日、12 月 21 日四个指定日期的表 1、表 2 与表 3 结果}。问题四在实时电价下分别按``仅 0:00 计划''与``0:00 计划＋三次调整''两种决策结构重算，并按同样的表式给出结果。全部策略依次保存到 \texttt{result1.xlsx}--\texttt{result4-3.xlsx}。
